\section{Введение}
Цель работы --- создать базу знаний об игровых персонажах, реализовать логические запросы на Prolog и перенести факты и отношения в онтологию, пригодную для работы в Protégé. Предметная область представляет собой небольшой вымышленный мир RPG: персонажи владеют оружием, обладают навыками, состоят в гильдиях и могут подходить для определённых заданий.

Практическая задача состоит в получении объяснимого ответа на вопрос: какой персонаж имеет оружие и навык, необходимый для задания? Например, для посещения башни требуется магия, поэтому при заданных фактах подходит Celia. Ответ определяется явно заданными условиями, а не статистическим прогнозом.

Работа охватывает обе части лабораторной 1: Prolog и онтологическое представление. Структура отчёта соответствует шаблону модуля 1. Диалоговая рекомендательная программа из лабораторной 2 относится к дальнейшему развитию и не входит в реализованный объём этой работы.

\section{Анализ требований}
\subsection{Функциональные требования}
База должна хранить типы объектов, роли персонажей, отношения владения, навыки и требования заданий. Необходимо поддерживать поиск с переменными, конъюнкцию, дизъюнкцию, отрицание и вывод по правилам. Результат подбора должен иметь проверяемое объяснение: предмет типа Weapon и совпадающий навык.

Для онтологии требуются классы, иерархия, свойства с областями определения и значений, индивиды и проверка выводимых знаний. Семантические различия Prolog и OWL должны быть описаны явно.

\begin{table}[ht]
\centering\small
\begin{tabularx}{\textwidth}{@{}Xll@{}}
\toprule
Требование задания & Реализация & Проверка\\
\midrule
Не менее 20 унарных фактов & 24 & Разбор исходника\\
От 10 до 15 бинарных фактов & 15 & Разбор исходника\\
От 5 до 7 правил & 7 & \code{knowledge.pl}\\
Разнообразные запросы & 15 тестов & SWI-Prolog\\
Комментарии к фактам и правилам & На русском языке & Исходный код\\
Перенос базы в онтологию & OWL и SWRL & Pellet, SPARQL\\
Сравнение двух представлений & 7 сопоставлений & Проверка множеств\\
\bottomrule
\end{tabularx}
\caption{Соответствие количественным и функциональным требованиям}
\end{table}

\subsection{Ограничения модели}
Данные являются учебными и вымышленными. Готовность к заданию означает только наличие оружия и подходящего навыка: уровень, здоровье, сложность, стоимость и вероятность победы не моделируются. Состояние Cursed рассматривается отдельно и не включено в правило готовности. Персонаж может иметь несколько ролей: несовместимость ролей не постулируется.

В базе хранятся шесть персонажей, три вида оружия, три задания, три навыка и две гильдии. Повторное использование \code{sword} означает владение видом оружия, а не одним уникальным физическим мечом. Имена являются стабильными идентификаторами объектов.

\section{Изучение основных концепций и инструментов}
\subsection{База знаний и логический вывод}
Факт \code{warrior(aria)} задаёт свойство объекта, а \code{owns(aria,sword)} --- отношение двух объектов. Правило задаёт условие вывода нового утверждения. Prolog использует унификацию и поиск с возвратом: переменная связывается с подходящим значением, а при запросе альтернатив перечисляются остальные решения.

В записи правила запятая обозначает «и», точка с запятой --- «или», оператор \code{\+} --- отрицание как неуспех доказательства. В \code{safe_character(X)} сначала выполняется \code{character(X)}, поэтому отрицание проверяется для конкретного известного персонажа. Перестановка этих целей при свободной переменной изменила бы поведение программы.

\subsection{Онтология и выбранные средства}
OWL описывает классы, свойства и ограничения, а механизм рассуждений выводит дополнительные типы и отношения. Protégé служит редактором и средством исследования такой модели. Онтология сохранена в RDF/XML и не требует внешних импортов. Для четырёх унарных правил используются определения классов OWL; для трёх правил, связывающих переменные через общий навык или гильдию, используются SWRL-правила.

Проверка Prolog выполняется в настоящем SWI-Prolog, поставляемом в пакете \code{swipl-wasm}, под Node.js. Скрипт \code{run_prolog.mjs} загружает исходник и сравнивает результаты с эталонами. Онтология строится через Owlready2 0.48, проверяется Pellet 2.3.1 под Java 17, а SPARQL-запросы выполняются RDFLib 7.6.0 по графу с материализованными выводами. SPARQL сам по себе не заменяет OWL/SWRL-рассуждение.

\section{Реализация системы искусственного интеллекта}
\subsection{Объекты и факты}
Унарные факты распределены следующим образом: 6 фактов Character, 6 фактов ролей, 3 факта Weapon, 3 факта Quest, 3 факта Skill, 2 факта Guild и 1 факт Cursed. Итого: 24 факта о 17 именованных объектах. Число фактов не совпадает с числом объектов, поскольку один персонаж имеет несколько свойств.

\begin{table}[ht]
\centering\small
\begin{tabularx}{\textwidth}{@{}lXXXX@{}}
\toprule
Персонаж & Роль & Оружие & Навык & Гильдия\\
\midrule
aria & warrior & sword & melee & sun\\
borin & warrior & sword & --- & moon\\
celia & mage & staff & magic & sun\\
darin & healer & staff & --- & ---\\
elin & rogue & bow & --- & ---\\
finn & ranger & bow & tracking & ---\\
\bottomrule
\end{tabularx}
\caption{Исходные сведения о персонажах; прочерк означает отсутствие факта}
\end{table}
Для Borin указан факт \code{cursed(borin)}. Задания требуют следующих навыков: \code{ruins} --- \code{melee}, \code{forest} --- \code{tracking}, \code{tower} --- \code{magic}. Бинарные факты включают 6 отношений владения, 3 навыка, 3 требования и 3 отношения членства: всего 15.

\subsection{Семь правил Prolog}
В листинге приведена полная логика вывода. Факты и русские комментарии находятся в приложенном файле \code{knowledge.pl}.
\begin{tcolorbox}[title={\texttt{knowledge.pl --- правила}},colback=white,colframe=black!70,breakable]
\begin{lstlisting}
combatant(X) :-
    character(X), (warrior(X); mage(X)).
scout(X) :-
    character(X), (rogue(X); ranger(X)).
armed(X) :-
    character(X), owns(X,W), weapon(W).
skilled_for(X,Q) :-
    character(X), quest(Q),
    has_skill(X,S), requires(Q,S).
ready_for(X,Q) :-
    armed(X), skilled_for(X,Q).
guildmate(X,Y) :-
    character(X), character(Y),
    member_of(X,G), member_of(Y,G), dif(X,Y).
safe_character(X) :-
    character(X), \+ cursed(X).
\end{lstlisting}
\end{tcolorbox}
Первые два правила классифицируют персонажей по альтернативным ролям. Третье проверяет наличие оружия. Четвёртое соединяет отношения по общей переменной навыка \code{S}; пятое объединяет результаты двух правил. Шестое находит разных членов одной гильдии и исключает пару объекта с самим собой. Седьмое реализует отрицание в границах известных персонажей.

\subsection{Перенос в онтологию}
Унарные предикаты перенесены в классы, факты --- в принадлежность индивидов этим классам. Бинарные предикаты перенесены в объектные свойства. Всего определены 15 классов, 7 объектных свойств и 17 предметных индивидов. Три из семи свойств предназначены для выводимых отношений.

\begin{table}[ht]
\centering\small
\begin{tabularx}{\textwidth}{@{}lXX@{}}
\toprule
Свойство & Domain & Range\\
\midrule
\code{owns} & Character & Weapon\\
\code{hasSkill} & Character & Skill\\
\code{requiresSkill} & Quest & Skill\\
\code{memberOf} & Character & Guild\\
\code{skilledFor} & Character & Quest\\
\code{readyFor} & Character & Quest\\
\code{guildmate} & Character & Character\\
\bottomrule
\end{tabularx}
\caption{Объектные свойства онтологии}
\end{table}
Классы Character, Weapon, Quest, Skill и Guild попарно несовместимы. Warrior, Mage, Healer, Rogue, Ranger и Cursed являются подклассами Character. Области Domain и Range используются для логического вывода типов, а не как процедура отклонения некорректных строк данных.

\begin{figure}[H]
\centering
\begin{tikzpicture}[n/.style={draw,rounded corners,align=center,font=\small,minimum height=8mm},>=Stealth]
\node[n] (c) at (0,0) {Character};
\node[n] (w) at (5,1.5) {Weapon};
\node[n] (s) at (5,-2) {Skill};
\node[n] (g) at (5,0) {Guild};
\node[n] (q) at (0,-2) {Quest};
\draw[->] (c) -- node[above,sloped,font=\small]{owns} (w);
\draw[->] (c) -- node[above,sloped,font=\small]{hasSkill} (s);
\draw[->] (c) -- node[above,font=\small]{memberOf} (g);
\draw[->] (q) -- node[below,sloped,font=\small]{requiresSkill} (s);
\draw[->] (c) -- node[left,font=\small]{readyFor} (q);
\end{tikzpicture}
\caption{Схема основных отношений модели; схема построена по онтологии, не является снимком интерфейса Protégé}
\end{figure}

Определения классов в синтаксисе Manchester имеют вид:
\begin{lstlisting}
Combatant EquivalentTo Character and (Warrior or Mage)
Scout EquivalentTo Character and (Rogue or Ranger)
Armed EquivalentTo Character and (owns some Weapon)
SafeCharacter EquivalentTo Character and not Cursed
\end{lstlisting}
Правила с общими переменными представлены в SWRL:
\begin{lstlisting}
Character(?x), Quest(?q), hasSkill(?x,?s),
requiresSkill(?q,?s) -> skilledFor(?x,?q)

Armed(?x), skilledFor(?x,?q) -> readyFor(?x,?q)

Character(?x), Character(?y), memberOf(?x,?g),
memberOf(?y,?g), differentFrom(?x,?y)
  -> guildmate(?x,?y)
\end{lstlisting}
Для всех предметных индивидов задано AllDifferent. Это необходимо, поскольку в OWL разные имена не гарантируют различие объектов; в Prolog различные атомы различаются непосредственно.

\subsection{Отрицание и границы эквивалентности}
В Prolog отсутствие доказательства \code{cursed(aria)} позволяет получить \code{safe_character(aria)}. В OWL при открытом мире отсутствие принадлежности Cursed не доказывает принадлежность дополнению этого класса. Поэтому для пяти непроклятых персонажей текущего снимка явно задан тип \code{not Cursed}. Это дополнительное предположение о полноте сведений о проклятии, а не автоматическое следствие отсутствия факта.

При добавлении нового персонажа без информации о проклятии Prolog сочтёт его безопасным, а OWL оставит классификацию неопределённой. При добавлении \code{cursed(aria)} необходимо пересмотреть отрицательное OWL-утверждение: сохранение обоих утверждений даст противоречие. Следовательно, равенство результатов относится к текущему снимку и оговорённым предположениям, а не ко всем возможным расширениям базы.

\subsection{Воспроизводимая проверка и работа в Protégé}
Сначала выполняется \code{node run_prolog.mjs}, затем \code{python build_ontology.py}. Второй скрипт извлекает исходные факты из Prolog, сохраняет \code{game_ontology.owl}, запускает Pellet и создаёт отдельный \code{results/game_inferred.owl}. После этого проверяются SPARQL-запросы и сопоставляются результаты семи правил с результатами SWI-Prolog.

Для демонстрации в Protégé следует открыть исходный OWL-файл через File / Open, посмотреть вкладки Classes, Object properties и Individuals. В дереве Character доступны роли и производные классы. Визуализацию иерархии можно открыть через OntoGraf при наличии этого представления. Для вывода SWRL-отношений требуется поддерживающий их reasoner, например Pellet; простой просмотр исходных RDF-троек такие отношения не вычисляет.

Материализованный файл можно открыть отдельно, чтобы увидеть уже проверенные отношения, например \code{aria readyFor ruins}. В DL Query выражение \code{Character and (readyFor value ruins)} выбирает Aria, а \code{Combatant and SafeCharacter} --- Aria и Celia при доступном выводе или загруженных материализованных типах. Автоматические результаты в этом отчёте получены через Pellet и RDFLib; ручная сессия и снимки экрана Protégé не использовались как подтверждение тестов.

\section{Оценка и интерпретация результатов}
\subsection{Запросы к базе знаний}
В таблице показаны фактически полученные результаты. Пустое множество означает отсутствие решений; для полностью конкретного запроса это соответствует \code{false}. Для положительного конкретного запроса указано \code{true}.
\small
\begin{longtable}{@{}p{0.49\textwidth}p{0.45\textwidth}@{}}
\toprule
Запрос Prolog & Результат\\
\midrule\endhead
\code{character(aria)} & true\\
\code{character(dragon)} & false\\
\code{warrior(X)} & aria, borin\\
\code{owns(X,bow)} & elin, finn\\
\code{combatant(X)} & aria, borin, celia\\
\code{scout(X)} & elin, finn\\
\code{armed(X)} & aria, borin, celia, darin, elin, finn\\
\code{skilled_for(X,tower)} & celia\\
\code{ready_for(X,Q)} & (aria, ruins), (celia, tower), (finn, forest)\\
\code{guildmate(X,Y)} & (aria, celia), (celia, aria)\\
\code{safe_character(X)} & aria, celia, darin, elin, finn\\
\code{ready_for(borin,ruins)} & false\\
\code{guildmate(aria,aria)} & false\\
\code{combatant(X), safe_character(X)} & aria, celia\\
\code{safe_character(dragon)} & false\\
\bottomrule
\caption{Результаты 15 тестов SWI-Prolog}
\end{longtable}
\normalsize
Контрпример Borin показывает, что наличия меча недостаточно: у него не задан навык melee. Elin и Darin также вооружены, но не имеют подходящих навыков в текущей базе. Пары членов гильдии симметричны вследствие тела правила, хотя симметричность свойства отдельно не объявлена. Проверка пары Aria с собой подтверждает действие ограничения различия.

\subsection{Запросы к онтологии и сравнение}
Пример SPARQL-запроса по материализованному графу:
\begin{lstlisting}
PREFIX : <https://example.org/rpg-lab1#>
SELECT ?x ?q WHERE { ?x :readyFor ?q }
\end{lstlisting}
Получены пары \code{aria--ruins}, \code{celia--tower} и \code{finn--forest}. Порядок выдачи без ORDER BY не гарантируется, поэтому при сопоставлении результаты нормализуются и сортируются.

Для класса безопасных персонажей используется положительный запрос к выведенному типу:
\begin{lstlisting}
PREFIX : <https://example.org/rpg-lab1#>
SELECT ?x WHERE { ?x a :SafeCharacter }
\end{lstlisting}
Подмена этого запроса фильтром NOT EXISTS по Cursed означала бы проверку отсутствия тройки в графе, а не доказательство OWL-отрицания. В реализации такая подмена не используется.

\begin{table}[ht]
\centering\small
\begin{tabularx}{\textwidth}{@{}Xrr@{}}
\toprule
Проверка & Количество & Успешно\\
\midrule
Запросы SWI-Prolog & 15 & 15\\
Запросы SPARQL после Pellet & 9 & 9\\
Сопоставления Prolog и OWL & 7 & 7\\
\bottomrule
\end{tabularx}
\caption{Итог автоматической проверки}
\end{table}
Pellet завершил рассуждение без обнаружения противоречия. Девять запросов охватывают четыре производных класса, три отношения и два отрицательных случая: отсутствие задания для Borin и отсутствие отношения guildmate к самому себе. JSON-протоколы содержат сами запросы, ожидаемые и фактические результаты. Проверки подтверждают текущий набор сценариев, но не являются доказательством корректности любых будущих изменений.

\subsection{Интерпретация и развитие}
Все количественные требования выполнены. Для учебной базы Prolog удобен компактностью правил и прямым интерактивным поиском. Онтология добавляет явную структуру типов, описание свойств и контроль логической согласованности, но требует осторожного обращения с отрицанием и различием индивидов.

На данном небольшом наборе не проводилось измерение производительности; утверждения о скорости или масштабируемости не делаются. Следующие шаги: добавить уровни и ограничения заданий, уточнить полноту сведений о состояниях персонажей, выдавать цепочку объяснения результата и реализовать диалоговый интерфейс лабораторной 2. При росте базы полезны отдельные проверки качества исходных данных и регрессионные тесты.

\section{Заключение}
Разработана прокомментированная база знаний из 24 унарных фактов, 15 бинарных фактов и 7 правил. Выполнены запросы разных типов, включая поиск с переменными, логические операции, составной вывод и отрицательные случаи. Подготовлена OWL/SWRL-онтология с иерархией классов, областями свойств и явными предположениями, необходимыми для сравнения с Prolog.

Получены три объяснимые пары персонажей и заданий; результаты семи правил согласуются в двух представлениях при фиксированном снимке знаний. Модель может служить основой для помощника игрока, проверки игровых данных или учебной системы поддержки принятия решений. Основной вывод работы: перенос знаний требует согласования семантики, а не только замены синтаксиса фактов.

\section*{Источники}
\addcontentsline{toc}{section}{Источники}
\begin{enumerate}
\item Задание курса AI System. \url{https://sunnysubmarines.notion.site/AI-System-a559a46cddc44363bdf27b77e10b7d85}.
\item Шаблон отчёта по модулю 1. \url{https://sunnysubmarines.notion.site/1-ea9a418d514b4baa8f7b4fc4c3ee176e}.
\item SWI-Prolog. Официальный сайт и документация. \url{https://www.swi-prolog.org/}.
\item Protégé. Официальный сайт. \url{https://protege.stanford.edu/}.
\item W3C. OWL 2 Structural Specification and Functional-Style Syntax. \url{https://www.w3.org/TR/owl-syntax/}.
\end{enumerate}

\appendix
\section{Файлы проекта и воспроизведение}
\begin{description}[style=nextline,leftmargin=0pt,labelwidth=0pt]
\item[\code{knowledge.pl}] Все факты, правила и комментарии на русском языке.
\item[\code{run_prolog.mjs}] Исполнение 15 запросов в SWI-Prolog и проверка ответов.
\item[\code{build_ontology.py}] Перенос фактов, OWL/SWRL-аксиомы, запуск Pellet, SPARQL и сопоставление результатов.
\item[\code{game_ontology.owl}] Исходная онтология для открытия в Protégé.
\item[\code{queries.sparql}] Девять отдельных запросов к графу с выводами.
\item[\code{results/}] JSON-протоколы и онтология с материализованными выводами.
\item[\code{report/main.tex, title.tex, 1.tex}] Исходный текст настоящего отчёта.
\end{description}
Команды выполняются из каталога Lab1. Требуются Node.js с npm, Python и Java. Зависимости устанавливаются из файлов проекта:
\begin{lstlisting}
npm ci
python -m pip install -r requirements.txt
node run_prolog.mjs
python build_ontology.py
\end{lstlisting}
Альтернативно Prolog можно запустить в настольном SWI-Prolog, загрузив \code{knowledge.pl}, и вводить запросы из таблицы вручную. Для сборки отчёта из каталога \code{report}:
\begin{lstlisting}
xelatex main.tex
xelatex main.tex
% or:
tectonic main.tex
\end{lstlisting}
Оформление соответствует предоставленному образцу: формат A4, кегль 11 pt, поля 2 см слева, справа и сверху и 3 см снизу, шрифты CMU Serif и CMU Typewriter Text, отдельные файлы титульного листа и основной части, оглавление и нумерованные листинги. Файлы шрифтов с лицензией включены в каталог \code{fonts}. Данные студента находятся в \code{title.tex}, имя преподавателя задаётся командой \code{TeacherName} в \code{main.tex}.
